\begin{ouverture}[courbe]{4}{Calcul littéral et identités remarquables}{Partie 1 · Nombres et calculs}
{Expressions · Distributivité · Identités remarquables · Preuves}
\objectif{Transformer une expression sans changer sa valeur.}
\objectif{Développer et factoriser avec les identités remarquables.}
\objectif{Démontrer une propriété générale avec une lettre.}
\prerequis{Produits, sommes, puissances et calculs avec les relatifs.}
\end{ouverture}
\pagestyle{otMain}\essentiel\fontsize{11}{15}\selectfont
\essSection{1}{Lire, simplifier et développer}
\begin{definition}Une expression littérale décrit un calcul avec des lettres. Sa nature est donnée par la dernière opération effectuée : $3x+2$ est une somme, $5(x+4)$ est un produit. Développer transforme un produit en somme ; factoriser transforme une somme en produit.\end{definition}
\begin{propriete}La distributivité donne $k(a+b)=ka+kb$. En particulier, $-(a-b)=-a+b$ : prendre l’opposé d’une expression change le signe de chacun de ses termes. Pour deux parenthèses :
\[(a+b)(c+d)=ac+ad+bc+bd.\]\end{propriete}
\begin{exemple}$(2x-3)(x+4)=2x^2+8x-3x-12=2x^2+5x-12$. Le terme $2x^2$ est un carré multiplié par 2, pas un terme en $x$.\end{exemple}
\begin{savoirfaire}[Simplifier un produit ou un rapport]Regrouper les facteurs de même nature : $(3a^2b)\times(2bc)=6a^2b^2c$. Dans un rapport, simplifier les facteurs communs non nuls : $\dfrac{6ab^2}{9a^2b}=\dfrac{2b}{3a}$.\end{savoirfaire}
\begin{attention}$3x+2x=5x$, mais $3x+2$ ne se réduit pas à $5x$. Un test numérique peut détecter une erreur ; il ne démontre pas une égalité pour tous les nombres.\end{attention}
\clearpage
\essSection{2}{Identités remarquables et factorisation}
\begin{propriete}[Identités remarquables]Pour tous nombres $a$ et $b$ :
\[(a+b)^2=a^2+2ab+b^2,\qquad (a-b)^2=a^2-2ab+b^2,\qquad (a-b)(a+b)=a^2-b^2.\]
Lues de gauche à droite, elles développent ; lues de droite à gauche, elles factorisent.\end{propriete}
\begin{exemple}
\[4x^2-25=(2x-5)(2x+5),\qquad x^2+10x+25=(x+5)^2,\qquad 99^2=(100-1)^2=10000-200+1=9801.\]
\end{exemple}
\begin{savoirfaire}[Factoriser]Chercher d’abord un facteur commun apparent : $6x^2-9x=3x(2x-3)$ et $3x(x-2)+5(x-2)=(x-2)(3x+5)$. Sinon, reconnaître une identité remarquable en identifiant $a$ et $b$.\end{savoirfaire}
\begin{savoirfaire}[Prouver avec un programme de calcul]Nommer le nombre choisi $x$, traduire chaque étape, puis transformer l’expression. Pour « ajouter 4, multiplier par 3, retirer le triple du nombre initial », on obtient $3(x+4)-3x=12$ : le résultat est indépendant de $x$.\end{savoirfaire}
\begin{attention}$(a+b)^2$ n’est généralement pas $a^2+b^2$. Il faut les deux produits croisés, dont la somme vaut $2ab$.\end{attention}
\clearpage\rubrique{Découverte}
\begin{blocLivre}{ex-4-1}\begin{exercice}[1]{Remplacer une lettre}
Calculer $A=2x^2-3x+1$ pour $x=-2$ et pour $x=3$.
\end{exercice}\end{blocLivre}
\begin{blocLivre}{ex-4-2}\begin{exercice}[1]{Somme ou produit ?}
Donner la nature de chaque expression : $A=3x+2$ ; $B=5(x+4)$ ; $C=(x-1)(x+3)$ ; $D=x^2-9$ ; $E=2x(x+1)+3$.
\end{exercice}\end{blocLivre}
\begin{blocLivre}{ex-4-3}\begin{exercice}[1]{Prendre l’opposé}
Supprimer les parenthèses et réduire $-(4x-7)+2(3-x)$.
\end{exercice}\end{blocLivre}
\begin{blocLivre}{ex-4-4}\begin{exercice}[1]{Double distributivité}
Développer et réduire $(3x+2)(x-5)$.
\end{exercice}\end{blocLivre}
\begin{blocLivre}{ex-4-5}\begin{exercice}[1]{Un facteur commun}
Factoriser $14x+21$ et $12x^2-8x$.
\end{exercice}\end{blocLivre}

\clearpage\rubrique{Entraînement}
\begin{blocLivre}{ex-4-6}\begin{exercice}[2]{Simplifier des produits et des rapports}
Simplifier, pour $a$, $b$, $c$ et $R$ non nuls : $(3a^2b)\times(2bc)$ ; $\dfrac{6ab^2}{9a^2b}$ ; $\dfrac{\frac43\pi R^3}{\pi R^2}$.
\end{exercice}\end{blocLivre}
\begin{blocLivre}{ex-4-7}\begin{exercice}[2]{Différences de carrés}
Factoriser $x^2-81$ et $25x^2-4$.
\end{exercice}\end{blocLivre}
\begin{blocLivre}{ex-4-8}\begin{exercice}[2]{Deux carrés de sommes}
Développer $(x+6)^2$ et $(2x-3)^2$.
\end{exercice}\end{blocLivre}
\begin{blocLivre}{ex-4-9}\begin{exercice}[2]{Reconnaître un carré}
Factoriser $x^2+10x+25$, puis $9x^2-12x+4$.
\end{exercice}\end{blocLivre}
\begin{blocLivre}{ex-4-10}\begin{exercice}[2]{Factoriser en deux temps}
Factoriser $3x(x-2)+5(x-2)$, puis $9x^2-36$.
\end{exercice}\end{blocLivre}

\clearpage\rubrique{Problèmes}
\begin{blocLivre}{ex-4-11}\begin{exercice}[3]{Deux programmes équivalents}
Programme A : ajouter 5 au nombre choisi, multiplier par 4, retirer 8. Programme B : multiplier le nombre choisi par 4, ajouter 12. Montrer qu’ils sont équivalents.
\end{exercice}\end{blocLivre}
\begin{blocLivre}{ex-4-12}\begin{exercice}[3]{Une bordure}
Un carré de côté $x$ cm est entouré d’une bordure de largeur 2 cm à l’extérieur. Exprimer l’aire de la bordure. Calculer pour $x=7$.
\label{fig-ex-4-12}\figureGeo{bordure}
\end{exercice}\end{blocLivre}
\begin{blocLivre}{ex-4-13}\begin{exercice}[3]{Des entiers et leur parité}
Montrer que le carré d’un entier impair est impair. Le carré d’un entier pair est-il toujours multiple de 4 ?
\end{exercice}\end{blocLivre}
\begin{blocLivre}{ex-4-14}\begin{exercice}[3]{Une égalité impossible ?}
Deux carrés ont pour côtés positifs $x+3$ et $x$. Leur différence d’aires peut-elle valoir 9 cm$^2$ ? Peut-elle valoir 39 cm$^2$ ?
\end{exercice}\end{blocLivre}
\begin{blocLivre}{ex-4-15}\begin{exercice}[3]{Une différence toujours positive}
Comparer, pour tout nombre $x$, $A=(x+1)^2$ et $B=x(x+2)$. Démontrer la réponse.
\end{exercice}\end{blocLivre}

\clearpage
